\documentclass[11pt]{article}
\usepackage[margin=1in]{geometry}
\usepackage{amsmath, amssymb}
\usepackage{booktabs}
\usepackage{setspace}
\usepackage{enumitem}

\setstretch{1.2}
\setlength{\parindent}{0pt}
\setlist[itemize]{itemsep=0.4em}
\setlist[enumerate]{itemsep=0.6em}

\begin{document}

\begin{center}
{\Large \textbf{Practice Problem Set 1}}\\
\vspace{1em}
{\large ECON 204 -- Contemporary Economic Issues}\\
\vspace{1em}
{\large Winter 2026}
\end{center}
\hrule
\vspace{1em}

%====================================================
\section*{Part I: Inflation, Cost of Living, and Monetary Policy}

\subsection*{Question 1: Measuring Inflation}

The Consumer Price Index (CPI) was 152.0 last year and 158.1 this year.

\begin{enumerate}
\item Calculate the inflation rate.
\item If a worker’s nominal wage increased by 3.5\%, did their real wage rise or fall?
\item Explain what your result implies about purchasing power.
\end{enumerate}

%----------------------------------------------------

\subsection*{Question 2: Cost of Living and Household Budgets}

A household spends:

\begin{itemize}
\item 30\% on housing
\item 20\% on food
\item 15\% on transportation
\item 10\% on utilities
\item 25\% on other goods
\end{itemize}

During the year:

\begin{itemize}
\item housing costs rise 8\%
\item food prices rise 5\%
\item transportation rises 4\%
\item utilities rise 6\%
\item other goods rise 2\%
\end{itemize}

\begin{enumerate}
\item Calculate the household-specific inflation rate.
\item Explain why this household’s inflation rate may differ from the national CPI.
\item Which households are most affected when housing costs rise faster than other prices?
\end{enumerate}

%----------------------------------------------------

\subsection*{Question 3: Real Interest Rates and Monetary Policy}

The nominal interest rate is 5.5\% and expected inflation is 2.8\%.

\begin{enumerate}
\item Calculate the real interest rate.
\item Explain how the real interest rate influences consumption and borrowing decisions.
\item If the Bank of Canada raises the policy rate, explain the channels through which inflation may decline.
\end{enumerate}

%----------------------------------------------------

\subsection*{Question 4: Inflation Expectations}

Suppose workers expect inflation to rise from 2\% to 5\%.

\begin{enumerate}
\item Explain how wage bargaining may change.
\item Explain how firms may respond.
\item Describe how expectations can create persistent inflation.
\end{enumerate}

%====================================================
\section*{Part II: Income Inequality}

\subsection*{Question 5: Measuring Inequality}

Five households earn the following annual incomes:

\begin{center}
\begin{tabular}{cc}
\toprule
Household & Income (\$) \\
\midrule
A & 20,000 \\
B & 30,000 \\
C & 45,000 \\
D & 70,000 \\
E & 135,000 \\
\bottomrule
\end{tabular}
\end{center}

\begin{enumerate}
\item Calculate the mean income.
\item Calculate the median income.
\item Explain why the mean may differ from the median.
\end{enumerate}

%----------------------------------------------------
\newpage
\subsection*{Question 6: Income Shares}

Using the data above:

\begin{enumerate}
\item Calculate total income.
\item Calculate the income share of the top household.
\item Calculate the income share of the bottom two households.
\item Interpret what these shares imply about inequality.
\end{enumerate}

%----------------------------------------------------

\subsection*{Question 7: Policy and Redistribution}

Suppose the government introduces a transfer of \$8,000 to each of the two lowest-income households, financed by a \$16,000 tax on the highest-income household.

\begin{enumerate}
\item Calculate new incomes.
\item Calculate the new mean income.
\item Has inequality increased or decreased? Explain.
\end{enumerate}

%----------------------------------------------------

\subsection*{Question 8: Poverty Measurement}

Assume the poverty threshold is \$28,000.

\begin{enumerate}
\item Which households are below the poverty line?
\item Calculate the poverty rate.
\item Explain one limitation of using a poverty line to measure economic hardship.
\end{enumerate}
\newpage
%====================================================
\section*{Part III: Housing Affordability and Real Estate Markets}

\subsection*{Question 9: Price-to-Income Ratio}

A household earns \$95,000 annually. The average home price is \$760,000.

\begin{enumerate}
\item Calculate the price-to-income ratio.
\item Interpret what this ratio implies about affordability.
\item Explain why this measure may not fully capture affordability.
\end{enumerate}

%----------------------------------------------------

\subsection*{Question 10: Mortgage Payments and Interest Rates}

A household borrows \$500,000 with a 25-year mortgage.

\begin{enumerate}
\item Estimate how monthly payments change if the mortgage rate rises from 2.5\% to 5.5\%.
\item Explain why rising interest rates can worsen affordability even if house prices stabilize.
\end{enumerate}

%----------------------------------------------------

\subsection*{Question 11: Rental Affordability}

A renter earns \$3,800 per month. Rent increases from \$1,200 to \$1,850.

\begin{enumerate}
\item Calculate the rent-to-income ratio before and after.
\item At what threshold is housing considered unaffordable?
\item Explain how rising rents contribute to inequality.
\end{enumerate}

%----------------------------------------------------

\subsection*{Question 12: Supply Elasticity and Prices}

Explain how each of the following affects housing affordability:

\begin{enumerate}
\item restrictive zoning regulations
\item increased immigration
\item new transit infrastructure
\item a decline in mortgage interest rates
\end{enumerate}
\newpage
%====================================================
\section*{Part IV: Integrated Policy Analysis}

\subsection*{Question 13: Inflation and Housing}

Explain how high inflation and rising interest rates can simultaneously:

\begin{itemize}
\item slow house price growth
\item increase mortgage payments
\item raise rental demand
\end{itemize}

%----------------------------------------------------

\subsection*{Question 14: Inequality and Housing Wealth}

Explain how rising house prices can increase wealth inequality between:

\begin{itemize}
\item homeowners and renters
\item older and younger households
\end{itemize}



\vspace{1em}
\hrule
\vspace{0.5em}
\textit{End of Practice Set}

\end{document}
